\subsection{Finite tests on the surviving positive-density packet}
\label{sec:fd-density-detection}

The common survival interval has an exact consequence for the finite
tests. Fix \(0<\eta<c_{\rm MT}\). For each sufficiently large \(T\),
choose the \(M\ge(c_{\rm MT}-\eta)M_T\) branches given by
Theorem~\ref{thm:two-thirds-survival}, and write them as
\(x_1(t),\ldots,x_M(t)\). Throughout
\[
 I_{T,\eta}=[-c_\eta/(\log T)^2,0]
\]
they are distinct simple real zeros in \((T,2T)\). Define
\[
 \mathcal F_t(w)=\frac{H_t(i\sqrt w)}{H_t(0)},\qquad
 P_{T,t}(w)=\prod_{j=1}^{M}\left(1+\frac{w}{x_j(t)^2}\right),
 \qquad \mathcal E_{T,t}(w)=\frac{\mathcal F_t(w)}{P_{T,t}(w)}.
\]
The first expression denotes its even-power series, so no square-root
branch is chosen. Every factor in the denominator divides the entire
numerator. Hence \(\mathcal E_{T,t}\) is real entire and has constant
coefficient one.

For any real series \(A(w)=\sum_{n\ge0}a_nw^n\), with \(a_0=1\)
and \(a_n=0\) for \(n<0\), define \(m_*(A)\) to be the least
order of a negative minor of \([a_{j-i}]_{i,j\ge0}\), and define
\(R_*(A)\) to be the least index window \(0,\ldots,R\) containing
such a minor. Define \(j_*(A)\) to be the least \(N\ge1\) for which
\[
 J_N(A;X)=\sum_{n=0}^{N}(N)_n a_nX^n
\]
has a nonreal zero. Each empty minimum is \(+\infty\). These
definitions also apply to the complementary quotient without imposing
an extra sign assumption on its coefficients.

\begin{corollary}[Detection on a positive-density surviving packet]
\label{cor:fd-density-detection}
For every \(t\in I_{T,\eta}\), the coefficient Toeplitz matrix of
\(P_{T,t}\) is totally nonnegative and every Jensen polynomial of
\(P_{T,t}\) has only real zeros. Thus
\[
 m_*(P_{T,t})=R_*(P_{T,t})=j_*(P_{T,t})=+\infty
\]
uniformly on this common interval. The exact factorization above gives
\[
 \begin{split}
 m_*(\mathcal F_t)&\ge m_*(\mathcal E_{T,t}),\\
 R_*(\mathcal F_t)&\ge R_*(\mathcal E_{T,t}),\\
 j_*(\mathcal F_t)&\ge j_*(\mathcal E_{T,t}).
 \end{split}
\]
All three comparisons use the same time and permit infinite values.
\end{corollary}
\begin{proof}
For \(c>0\), the Toeplitz matrix of \(1+cw\) has diagonal entries
one, entries \(c\) immediately above the diagonal, and all other
entries zero. In a minor with increasing row and column lists, a
nonzero determinant term cannot contain an inversion. Indeed, for
rows \(i_p<i_q\), nonzero entries require their assigned columns to
satisfy \(j_p\le i_p+1\le i_q\le j_q\). An inverted assignment
would contradict this inequality. Thus the only possible nonzero
term has the positive identity-permutation sign, proving total
nonnegativity.

Toeplitz convolution is exact multiplication:
\[
 T(AB)=T(A)T(B).
\]
For each fixed minor the intermediate indices range over a finite
set, because the matrices are upper triangular. Cauchy--Binet
therefore proves total nonnegativity for the finite product
\(P_{T,t}\). More generally, when every minor of \(T(B)\) of orders
at most \(r\) is nonnegative, the same is true of \(T(P_{T,t}B)\).
This proves the comparison for \(m_*\). For the finite window,
\[
 T_R(P_{T,t}B)=T_R(P_{T,t})T_R(B),
\]
where \(T_R\) uses precisely the indices \(0,\ldots,R\).
Applying the same finite Cauchy--Binet formula proves the comparison
for \(R_*\), with no change in the window.

For the Jensen test, reverse the polynomial and write \(D=d/dx\):
\[
 Q_N(A;x)=x^N J_N(A;1/x)=A(D)x^N.
\]
Only the first \(N+1\) terms of the formal differential operator
act nontrivially. Coefficient convolution now gives the exact map
\[
 Q_N(P_{T,t}B;x)
   =P_{T,t}(D)Q_N(B;x)
   =\prod_{j=1}^{M}\left(1+\frac{D}{x_j(t)^2}\right)Q_N(B;x).
\]
Each operator \(Q\mapsto Q+cQ'\), for \(c>0\), preserves real
roots. For a polynomial with distinct real roots this follows from
\(Q'/Q=\sum_r(x-r)^{-1}\), which decreases strictly from
\(+\infty\) to \(-\infty\) between consecutive roots and gives
one additional solution of \(Q'/Q=-1/c\) to the left of all roots.
Polynomials with repeated real roots follow by perturbing their
roots to distinct real ones and taking the coefficient limit;
real-rootedness is closed under that limit. Constants are fixed.
Starting from \(x^N\) therefore proves the assertion for the packet.
Starting from a real-rooted \(Q_N(B;x)\) proves it for the product.
Polynomial reversal preserves reality of nonzero roots and merely
adds or removes roots at zero; its leading coefficient here is one.
Thus \(J_N(B)\) real-rooted implies \(J_N(P_{T,t}B)\) real-rooted,
which proves the comparison for \(j_*\).
\end{proof}

This corollary transports the positive-density survival count to
infinite local detecting orders on one common interval. Its relation
to the global tests is the displayed product map: the complementary
quotient controls the lower bound for each global detecting index.
The question of divergent scaled global orders still asks for the
corresponding estimates for the full complementary factor as the
height and the backward time vary.
